\section{Research status and the concrete advance}
\label{sec:status}

\subsection{What is being continued}

The maintained repository already contains a substantial exact recognizer:
structural and braid certificates, polynomial one-sided obstructions,
width-sensitive Jones and Khovanov computations, compressed group operations,
and an optional complete normal-surface backend. Its synthesis article had
reached 210 pages before this work. Repeating the existing width-conditional
recurrences, or adding another invariant with no effect on the surviving
search space, would not close its main theoretical gap.

The public checkout and preceding relevant revision are pinned as follows
\cite{proveit}:
\begin{center}\small
\begin{tabular}{@{}ll@{}}
Checkout & \baseline\\
Subtree revision & \subtreebaseline\\
Shared \code{fast/} tree & \texttt{e7222244bb0d7f35cf7f34f83228b350240e4d13}
\end{tabular}
\end{center}
The preceding subtree revision is entitled ``Document sparse endpoint proofs
and controlled overlap measurements.'' Both references have the same
maintained \code{fast/} tree. The patch is prepared against the full checkout
reference. It does not depend on remembering an earlier conversational archive.

The latest compressed-overlap implementation already had two certified
shortcuts. One recognized a common cyclic pattern with period at most 32.
The other counted occurrences of an endpoint letter, saturating at 65, and
individually verified at most 64 forced candidates. Their exact fallbacks were
complete, but long-period blocks with many repeated markers could still
exhaust the allotted work. We remove those two numerical restrictions on a
structurally certified class. The important change is from enumerating a
bounded number of candidate lengths to proving a formula for all of them.

Two geometric interfaces are developed alongside that change. The existing
surface-cover kernel describes a particular affine action on a cyclic sheet
set. It does not solve gluing equivalence for an assembly of regular covers.
The new assembly kernel has a deliberately different, checked input model.
The existing Regina wrapper exports an external-engine verdict and metadata;
it does not independently certify a supplied normal disk. The new verifier
checks such a disk in a supplied finite triangulation. These are exact
additional capabilities, with their premises visible at the interface.

\subsection{What the current literature actually supplies}

Lackenby's July 2026 preprint gives a new algorithm based on essential
hierarchies and boundary patterns \cite{lackenby2026}. Section 9 explicitly
distinguishes a bound on the number of steps from a running-time estimate for
all steps. Its later sections develop polynomial verification and compressed
decomposition when suitable surfaces are supplied. In particular, supplying
a normal surface in binary coordinates and cutting efficiently along it are
not equivalent to finding the required surfaces or bounding the complete
recognition computation. The repository correctly preserves this distinction.

Compressed-word equality and several fully compressed matching operations
already have polynomial algorithms in grammar size
\cite{schleimer,lifshits,matsubara}. The present endpoint construction is a
more economical exact route for a broad promised class of overlap queries.
We do not claim that it changes a previously exponential general compressed
matching problem into a polynomial one.

Likewise, polynomial normal-surface verification is classical. The
extreme-ray approach is part of the Hass--Lagarias--Pippenger framework
\cite{hlp}; the more general orbit-counting machinery of
Agol--Hass--Thurston handles surfaces beyond this restricted certificate
class \cite{aht}. Our contribution is a small independently inspectable
implementation, an explicit modular rank certificate, a parity test for the
boundary, and a reproducible comparison that does not require expanding a
normal disk. General normal-surface discovery remains a separate problem.

\subsection{Results and their precise scope}

\begin{longtable}{@{}>{\raggedright\arraybackslash}p{.25\textwidth}>{\raggedright\arraybackslash}p{.38\textwidth}>{\raggedright\arraybackslash}p{.30\textwidth}@{}}
\toprule
Result & Exact guarantee & Remaining premise or limitation\\
\midrule
Endpoint progressions & All suffix--prefix overlaps are returned from a
certified short-anchor occurrence progression; both its step and count are
binary integers with no fixed cap. & Failure of the structural certificate
uses the complete matcher; the inherited extension queries can still be
expensive.\\
Regular-cover gluing & All framed isomorphisms are encoded by at most two
root congruence classes per connected assembly graph. & Each block is a
connected regular cover with checked compatible whole-boundary attachments;
the blocks are already supplied.\\
Gluing state count & Exact orbit counts, including $m^\beta$ for cyclic
covers, quantify states left after relabelling. & This is a lower bound for
enumerating these framed classes, not for arbitrary unknot algorithms.\\
Normal disk certificate & A supplied finite triangulation and binary normal
vector are checked to contain an essential compressing disk. & The vector
must satisfy the extreme-ray condition; correspondence with a given PD
diagram is not asserted.\\
\bottomrule
\end{longtable}

The distinction between a complete recognizer and a complete local operation
is central. The word shortcut participates in an already existing optional
compressed group search. The gluing and normal-disk kernels are new explicit
APIs. They do not silently introduce new positive or negative knot verdicts
into the pipeline.

\section{Encodings, cost measures, and proof obligations}
\label{sec:models}

The number $n$ denotes crossings in a validated classical knot diagram,
not expanded word length or the size of an unencoded geometric construction.
We use the specific target $n^{O(\log n)}=2^{O((\log n)^2)}$. Some authors use
``quasi-polynomial'' for the larger union of bounds
$2^{(\log n)^{O(1)}}$; any use of that broader convention must be stated.

An SLP is an acyclic binary grammar whose rules are terminals or
concatenations of earlier rules. Let $g$ be the number of reachable rules,
$H$ its height, $L$ the maximum represented length, and
$B=\lceil\log_2(L+1)\rceil$. Lengths, positions, and counts are binary
integers. Signed generator names have their own bit lengths and are compared
exactly. A grammar may contain two copies of the same child in one rule;
its represented occurrence count must therefore respect multiplicity even
when the grammar nodes themselves are shared.

For surface assemblies, $m$ is the cyclic parameter, and
$\ell=\lceil\log_2(m+1)\rceil$. The regular dihedral group has order $2m$;
it is not the generally nonregular action of a dihedral subgroup on $m$
residues. Let $v,e,r$ denote the numbers of blocks, attachments, and
explicit markings, with the total listed local monodromy data included in
the input size. For a connected assembly graph, $\beta=e-v+1$.

For the normal-disk verifier, $t$ is the number of tetrahedra, $x$ the
normal coordinate vector, $s=|\supp x|\leq7t$, and $B_x$ the largest
coordinate bit length. A tetrahedron has four triangle and three
quadrilateral types. Face maps are explicit vertex permutations, not an
unverified topology label or a hash of one.

We distinguish three measurements throughout. A mathematical bit bound counts
integer operations according to operand length. The implementation's work
counter charges selected operations cooperatively; it is a resource policy,
not an exact instruction count. Wall-clock measurements are empirical and
include the scope specified by each benchmark. A resource-limited computation
has no recorded completion time. It cannot serve as the denominator of a
completion-speedup claim.

\begin{proposition}[Local bounds do not compose without state bounds]
Suppose each operation on an encoded state of size $S$ costs $S^c$ for a
fixed $c$. If a computation performs $R(n)$ charged state-processing
operations, including repeated visits, and every participating state has
size at most $S(n)$, its operation cost is bounded by
$R(n)S(n)^c$, before adding input conversion and output verification.
To obtain the target $n^{O(\log n)}$ from this estimate, it suffices to prove both
$\log R(n)=O((\log n)^2)$ and
$\log S(n)=O((\log n)^2)$, with the other work similarly controlled.
\end{proposition}
\begin{proof}
Sum the local cost over the actual operations. Taking logarithms yields
$\log R+c\log S$. The statement requires bounds for the states actually
processed and the number of processing operations; the existence of a
smaller representation, or a bound only on distinct states, is insufficient.
\end{proof}

This elementary accounting is included to locate the obligation, not as a
new recognition theorem. The gluing orbit count below gives a concrete
reason that an efficient quotient computation can leave too many distinct
states for this argument to meet the target.
